\section{Advisor Search and Maintenance Constraint}

Search is a client of the native what-if substrate, not the paper's primary
technical novelty.  The product obtains a baseline with
\texttt{evaluate\_design} and evaluates counterfactual configurations through
\texttt{evaluate\_move}; PostgreSQL remains responsible for native
applicability, payload consumption, and cardinality estimation.  Search is
therefore replaceable and is not a new optimization algorithm.  The current
product uses a deterministic contextual ADD-only construction because it is a
small, reproducible client with which to exercise the evaluator under a
maintenance budget.

\subsection{Search as a substrate consumer}

For a current design $Y$, an ADD move for an unselected candidate $s$ asks the
substrate to evaluate $Y\cup\{s\}$ in the context of the active design.  The
search client does not replace the native evaluator with a utility model: it
enumerates candidate moves, supplies them to the evaluator, and compares the
returned values of $L(Y;R)$.  This separation means that another client could
use the same candidate repository, ordered activation, and evaluator without
adopting this search policy.

\subsection{Deterministic contextual ADD search}

The release-qualified workflow sets the \texttt{add\_only} field of
\texttt{SearchConfig} to true.  Its product algorithm is a
deterministic contextual ADD-only best-improvement (greedy) construction.  At
each round it enumerates every candidate not in $Y$, rejects a candidate whose
counterfactual design would violate $C(Y')\le B$, applies exact lower-bound
pruning where available, and natively evaluates the remaining moves.  It
accepts the best strictly improving move and repeats from the resulting
design.  A tie is resolved by the lower maintenance cost and then by the
candidate's deterministic rank key; for ADD moves this key is precedence rank
followed by candidate ID.  Candidate enumeration and activation use the same
fixed precedence discipline, so equal inputs produce the same trajectory.

The contextual marginal change for a candidate is
\[
 \Delta(s\mid Y,R)=L(Y;R)-L(Y\cup\{s\};R),
\]
not a fixed singleton score.  The native estimate for $Y\cup\{s\}$ can
depend on the currently active statistics, on precedence, and on other
selected candidates.  Contextual greedy/local physical-design search is
established prior art; the role of the algorithm here is to exercise the
native what-if substrate under a scoped maintenance constraint.  The library
also contains an ADD/DROP/SWAP development mode, but that mode is not the
release-qualified advisor path evaluated as the product workflow.

\subsection{Maintenance constraint}

The search uses the notation of Section~\ref{sec:problem}.  Each candidate
definition has a cost $c(s)$, the design cost is
$C(Y)=\sum_{s\in Y}c(s)$, and a move is feasible only when its endpoint
satisfies $C(Y')\le B$.  The current model is an empirical mechanism-count
model for the supported arity-two MCV/FD catalog.  Its parameters are fixed
before search, validated against the externally supplied target $T$, and
loaded as part of the run configuration; they are not learned or refit while
the search proceeds.  The model is target- and environment-specific.  It is a
reproducible first-order maintenance resource constraint, not a universal
PostgreSQL maintenance law or a causal per-candidate \texttt{ANALYZE} model.

Cost belongs to the candidate definition, rather than to the realized
payload.  Thus a candidate with either \texttt{PRESENT} or
\texttt{ABSENT\_NATIVE} state has the same definition-level cost.  The budget
check is performed before native evaluation whenever the implementation can
reject the move directly, avoiding work for infeasible endpoints while
preserving the design semantics from Section~\ref{sec:problem}.

\subsection{Contextual utility}

The search evaluates utility in the current context.  A candidate that is
weak or even harmful as a singleton can improve a design after another
candidate has changed native applicability, clause consumption, or MCV/FD
composition.  Conversely, a positive singleton result need not determine the
best move later in the trajectory.  The evaluator therefore computes
$\Delta(s\mid Y,R)$ for the current $Y$ rather than caching one design-
independent score.  This is an application of established statistics
interaction behavior, not a new interaction theorem.

The release-qualified \texttt{advise\_fixed\_t} path uses the full product
candidate catalog as its search-visible set.  Singleton profiles and explicit
candidate-set filters are available as separately versioned development or
artifact interfaces, but the core workflow does not claim that a negative
singleton utility is a generally safe pruning rule.  If a restricted visible
set is supplied by another workflow, that restriction is an explicit input
and must not be confused with a theorem about the omitted candidates.

\subsection{Termination and nonclaims}

The high-level loop is:
\[
\begin{aligned}
Y&\leftarrow\varnothing;\\
&\textbf{repeat:}\\
&\quad\text{enumerate feasible ADD moves;}\\
&\quad\text{prune only with the exact bound;}\\
&\quad\text{evaluate remaining moves natively;}\\
&\quad\text{choose the deterministic best strict improvement;}\\
&\textbf{until none exists.}
\end{aligned}
\]
The returned search result contains the selected design $Y$, its objective
$L(Y;R)$, its maintenance cost $C(Y)$, the budget $B$, and deterministic
trajectory and operation metadata.  The product termination label is
\texttt{add-local-optimum}: no feasible single ADD move gives a strict
objective improvement under the current budget and frozen realization.

The evaluator's conservative incidence index supplies the reuse boundary for
each candidate move.  The endpoint candidate determines an affected query set
$A$; PostgreSQL replans $A$ under the counterfactual ordered activation, while
the evaluator reuses the contributions for $Q\setminus A$.  False positives
in $A$ cost additional planner work, whereas a false negative would violate
correctness.  Incrementality therefore stops at query granularity, as
described in Section~5.

Section~5 also gives the safe objective lower bound used
before native evaluation.  A move is skipped only when the bound certifies no
strict improvement or that the move cannot beat the current best move.  The
pruning is exact for the stated q-error objective and conservative incidence
contract, but it is a secondary search-client optimization rather than a new
algorithmic contribution.

The product path makes no claim of search novelty, global optimality, or an
approximation guarantee.  Its ``local optimum'' qualifier is specifically
ADD-local; it is not a DROP/SWAP local optimum.  The maintenance model is
scoped to its fixed target and calibration environment, singleton utility is
not a universally safe pruning rule, and query reuse does not extend to
planner-internal state.
